\chapter*{常见缩写和术语先导}
\addcontentsline{toc}{chapter}{常见缩写和术语先导}

本书会从零开始讲 C++，但工程资料里很多词会以缩写或英文术语出现。先把这些词放在同一页，是为了让你读第一遍时不用猜字母含义，也不用把工具名、状态名和语言概念混在一起。

\begin{description}
  \item[ABI] Application Binary Interface，应用二进制接口。它规定函数调用、对象布局、符号和链接边界如何在二进制层面配合。
  \item[API] Application Programming Interface，应用程序接口。API 是源码层面的调用合同；ABI 是编译后二进制层面的合同。
  \item[ID] Identifier，身份或编号。示例里的用户 ID、玩家 ID 都是定位业务对象的键，不是“随便一个整数”。
  \item[C++14/C++20/C++23] C++ 标准版本名。C++14 引入泛型 lambda 等能力；C++20 是本书默认实现边界；C++23 更新了部分库能力，但你的编译器和刷题平台未必已经完整支持。
  \item[RAII] Resource Acquisition Is Initialization，资源获取即初始化。对象构造时拿资源，析构时释放资源，是 C++ 管理文件、锁、内存和线程句柄的核心习惯。
  \item[UB] Undefined Behavior，未定义行为。程序一旦触发 UB，标准不再规定结果；它可能看似正常，也可能被优化器改成完全不同的行为。
  \item[I/O/EOF] I/O 是 Input/Output，输入输出；EOF 是 End Of File，文件或输入流已经没有更多数据。读循环里要区分“读到数据”“读完”和“出错”。
  \item[ASCII/Unicode/UTF-8] ASCII 是早期单字节字符集合；Unicode 是统一字符集合；UTF-8 是 Unicode 的常见字节编码。C++ 的 \code{std::string} 保存字节，不自动理解用户眼里的“一个字”。
  \item[IDE/GUI] IDE 是集成开发环境，例如带编辑、构建、调试功能的工具；GUI 是图形用户界面。本书主要用命令行解释底层动作，不要求先依赖 IDE。
  \item[CI] Continuous Integration，持续集成。每次提交后自动构建、测试和检查，防止问题积累到发布前才暴露。
  \item[CLI] Command-Line Interface，命令行接口。综合项目里的配置工具、日志工具和文件索引器都会把功能做成 CLI。
  \item[CMake] C++ 工程常用构建系统。它把源文件、库、可执行文件、测试和编译选项组织成可重复构建的 target。
  \item[GNU] GNU 是自由软件工具链生态名，\code{g++} 是 GNU C++ 编译器驱动。它不是 C++ 语言本身，而是一套具体工具实现。
  \item[HTTP] HyperText Transfer Protocol，超文本传输协议。架构章节提到 HTTP 时，通常表示外部服务接口或网络适配器边界。
  \item[\code{std::}] C++ 标准库命名空间前缀。\code{std::vector}、\code{std::optional}、\code{std::span}、\code{std::string_view}、\code{std::mutex}、\code{std::future} 这类名字都来自标准库；先理解它们表达的所有权、视图、同步或异步语义，再记函数名。
  \item[\code{size\_t/int64\_t/uint64\_t}] 常见整数类型。\code{size_t} 用于大小和下标；\code{int64_t} 是 64 位有符号整数；\code{uint64_t} 是 64 位无符号整数。选择类型时看值域、是否允许负数和是否要和容器大小交互。
  \item[target/PUBLIC/PRIVATE/INTERFACE] target 是 CMake 中的构建目标；PUBLIC 依赖会传给使用者；PRIVATE 只影响当前 target；INTERFACE 只传播用法要求、不编译自身源码。
  \item[\code{DCMAKE\_BUILD\_TYPE/DCMAKE\_CXX\_FLAGS}] 常见 CMake 配置变量。前者选择 Debug、Release 等构建类型；后者追加 C++ 编译选项。它们会直接改变优化、断言、调试信息和诊断结果。
  \item[ASan/UBSan/TSan] 三类 sanitizer。ASan 查越界和 use-after-free；UBSan 查未定义行为；TSan 查数据竞争。它们服务正确性诊断，不用于发布性能数字。
  \item[TLS] Thread-Local Storage，线程局部存储。每个线程各自拥有一份变量实例，常用于线程内缓存或上下文。
  \item[JSON] JavaScript Object Notation，一种文本数据格式。C++ 项目常用它保存配置、测试输入或报告。
  \item[CPU] Central Processing Unit，中央处理器。本书讲 CPU 时，通常关心它如何执行你的 C++ 程序，而不是只关心硬件品牌。
  \item[O0/O2/DNDEBUG] 常见编译选项。O0 表示基本不优化，方便调试；O2 表示开启常用优化；DNDEBUG 会关闭依赖 \code{assert} 的调试检查。
  \item[LRU] Least Recently Used，最近最少使用淘汰。项目章会用它训练对象所有权、哈希表、链表和边界测试。
  \item[FIFO/LIFO] FIFO 是 First In First Out，先进先出，典型结构是队列；LIFO 是 Last In First Out，后进先出，典型结构是栈。
  \item[constexpr/type traits/concepts] C++ 模板和编译期编程常见词。constexpr 表示可编译期求值；type traits 查询类型属性；concepts 约束模板参数能力。
  \item[SFINAE] Substitution Failure Is Not An Error，模板替换失败不算硬错误的规则。它是早期模板约束技巧；C++20 之后很多场景更适合用 concepts 写清“这个模板要求参数具备什么能力”。
  \item[lambda] 局部定义的匿名可调用对象。可以把它理解成编译器帮你生成一个带 \code{operator()} 的小类；捕获列表就是这个小类里保存的成员。
  \item[span/string\_view] span 是一段连续对象的非拥有视图，通常保存指针和长度；string\_view 是字符串字节区间的非拥有视图。它们不延长底层对象生命周期。
  \item[future/worker] future 表示“稍后可取得的结果或异常”；worker 表示执行任务的线程或执行单元。异步代码里要分清提交任务、执行任务和取得结果这三步。
  \item[config/hash/token/padding] config 是配置；hash 是摘要或哈希值；token 是解析出的词元或文本片段；padding 是为满足对齐、协议或格式而填充的无业务字节。
  \item[\code{CPP/HPP}] 常见文件后缀。\code{.cpp} 通常放实现代码，\code{.hpp} 通常放 C++ 头文件；后缀只是约定，真正的边界由 include、target 和接口设计决定。
  \item[INFO/WARN/ERROR/UNKNOWN] 常见日志级别或状态名。INFO 表示普通信息，WARN 表示可恢复风险，ERROR 表示失败，UNKNOWN 表示当前没有足够信息分类。
  \item[命名常量] \code{PORT_MIN}、\code{PORT_MAX}、\code{RETRY_MAX_ATTEMPTS}、\code{USERNAME_MAX_LENGTH}、\code{LOGGER_MAX_QUEUE_SIZE} 这类全大写名字通常是业务或工程边界常量。它们不是新语法，而是把端口范围、重试次数、用户名长度和队列容量写成可审查的名字。
  \item[MB] Megabyte，兆字节量级。性能章节出现 MB 时，重点是数据规模和内存流量，不是 C++ 语言关键字。
  \item[sanitizer/benchmark] sanitizer 是编译器插桩诊断工具族，用来找越界、未定义行为或数据竞争；benchmark 是可重复的性能测量，不等于随手跑一次计时。
  \item[智能指针] \code{unique\_ptr} 表示独占所有权，不能复制，只能移动；\code{shared\_ptr} 表示引用计数共享所有权；\code{weak\_ptr} 只观察 shared 对象，不延长生命周期。读到这些名字时，先判断“谁负责释放对象、谁只是借用对象”。
  \item[常见异常类型] \code{runtime\_error} 表示运行时错误；\code{error\_code} 表示可比较、可传递的错误码。异常适合直接中断当前路径，错误码适合显式向调用者返回失败原因；二者都要配合清楚的接口合同。
  \item[C 文件接口] \code{FILE} 是 C 标准库里的文件流对象类型；C++ 项目里看到它，通常是为了演示“外部资源需要关闭”这个 RAII 问题。不要把 FILE 当成文件名，它是一个需要 \code{fclose} 或包装类管理的资源句柄。
  \item[低层类型转换] \code{reinterpret\_cast} 会把同一批位按另一个类型解释，常用于对象布局、字节查看或系统接口边界。它比 \code{static\_cast} 更危险，不能拿来“修复编译错误”；使用前必须说明对齐、生命周期和别名规则。
  \item[同步原语] \code{lock\_guard} 是作用域锁，构造时加锁、析构时解锁；\code{unique\_lock} 比 lock\_guard 更灵活，能延迟加锁、手动解锁并配合条件变量；\code{condition\_variable} 用于线程等待某个条件变化；\code{notify\_one/notify\_all} 分别唤醒一个或全部等待线程。
  \item[协程关键词] \code{co\_await} 表示当前协程可能暂停，等待某个 awaitable 对象完成；\code{await\_resume} 是恢复后取结果的接口。先把协程理解成“能把执行状态保存起来、稍后继续”的函数，再看语法糖。
  \item[编译期检查] \code{static\_assert} 在编译期验证条件；\code{has\_value} 常见于 \code{optional}，表示里面是否真的有值。它们都在减少“运行到一半才发现前提不成立”的情况。
  \item[CMake 命令] \code{add\_executable} 定义可执行程序 target；\code{target\_link\_libraries} 声明 target 依赖哪些库；\code{cmake -S . -B build} 通常表示从源码目录生成一个独立构建目录。读 CMake 时要跟踪 target，而不是只看文件名。
  \item[CMake 安装词] \code{TARGETS}、\code{RUNTIME}、\code{LIBRARY}、\code{ARCHIVE}、\code{DESTINATION}、\code{DIRECTORY} 是 CMake 安装和目标声明里的常见关键词。RUNTIME 多指可执行程序输出，LIBRARY 指共享库，ARCHIVE 指静态库或导入库，DESTINATION 指安装目录。
  \item[项目示例类型] \code{ScoreSummary}、\code{UserRecord}、\code{ConfigError}、\code{FileHandle}、\code{LruCache}、\code{AsyncLogger} 这类 CamelCase 名字是本书项目对象。它们不是标准库名，而是把业务角色写进类型名：摘要、用户记录、配置错误、文件句柄、缓存、异步日志器。
  \item[函数名读法] \code{read\_scores}、\code{parse\_port\_text}、\code{find\_by\_id}、\code{query\_database} 这类 snake\_case 名字通常按“动词 + 对象”读。先问它修改谁、返回谁、失败如何表达，再看函数内部代码。
\end{description}
